\chapter*{Abstract}

Interactive musical prediction systems are a significant subfield of New Interfaces for Musical Expression (NIME). These systems can perform collaboratively in a call-and-response manner with a user, with a neural network taking over from a human performer in real-time and predicting musical continuations. For musical predictions, there is never one correct answer for what the next note should be, and so neural networks usually rely on random sampling from a distribution of note options to generate a prediction. However, this random sampling is entirely arbitrary, and does not take into account how well each potential note fits into the context of the musical piece being played. This can lead to a lack of musical coherence in generated predictions. In contrast, older generative music techniques often used rule-based approaches, which were excellent at producing coherent music which adheres to musical rules, but lacked the expressiveness and creativity of modern machine learning based systems. This thesis introduces a novel musical prediction architecture which combines the strengths of both approaches, by using repeated neural network predictions to build a Monte Carlo search tree of viable musical continuations, and guiding the search with heuristics inspired by rules-based techniques. The evaluation measures improvements in prediction accuracy of 67.5\% and 145.2\% for two datasets, and qualitative analysis of improvisations performed with the system showed that this increase did not come at the cost of creativity, with the expressiveness and playability of the system being significantly improved by using this search based prediction architecture. This thesis represents a significant advancement for musical prediction systems and human-computer interaction more broadly, acting as a proof of concept for the prediction architecture developed in this thesis, and leading to future work possibilities for many domains where search based and machine learning based methods could be combined.